% SCP10, paper_v3.tex:2904--2923; placement and covariance only.
% Lean elaboration and current-head hosted acceptance remain pending.
\paragraph{Coherent color changes on the periodic lattice.}
Let $V=(\mathbb Z/W\mathbb Z)\times(\mathbb Z/H\mathbb Z)$ with $W,H\geq3$, and let $E$ be the edges of its nearest-neighbor simple graph.
Write $N_v,E_v,S_v,W_v$ for the north, east, south and west edges incident to $v$; in particular, $S_v=N_{v-(0,1)}$ and $W_v=E_{v-(1,0)}$.
The four physical spins at each vertex have the clockwise order of Definition~\ref{def:peps_quantum_double_k_tensor}.
All coordinate additions below are periodic, including at either seam.

\begin{definition}[Bond color action]
    \label{def:peps_qd_lattice_gauge}
    \lean{TNLean.PEPS.QuantumDoubleKLatticeConfig,
        TNLean.PEPS.quantumDoubleKLatticeGauge,
        TNLean.PEPS.quantumDoubleKLatticeGaugePermutation,
        TNLean.PEPS.quantumDoubleKLatticeBondPermutation}
    \uses{def:peps_quantum_double_k_tensor}
    On the full physical configuration space $\mathcal X=(G^4)^V$, let $u\in G^E$ act by
    \begin{align}
        (\Gamma_u x)_v
        & =\bigl(u_{N_v}a_vu_{E_v}^{-1},\,
        u_{E_v}b_vu_{S_v}^{-1},\,
        u_{S_v}c_vu_{W_v}^{-1},\,
        u_{W_v}d_vu_{N_v}^{-1}\bigr),
        \label{eq:qd_lattice_gauge}
    \end{align}
    where $x_v=(a_v,b_v,c_v,d_v)$.
    For $e\in E$ and $s\in G$, set $V_e(s)=\Gamma_{\delta_e(s)}$, where $\delta_e(s)$ equals $s$ at $e$ and $1$ elsewhere.
    This is the coherent bond action of \cite{Schuch2010PEPS}, source lines 2918--2923, with every physical register retained.
\end{definition}

\begin{lemma}[Composition and distinct-bond commutation]
    \label{lem:peps_qd_lattice_gauge_action}
    \lean{TNLean.PEPS.quantumDoubleKLatticeGauge_one,
        TNLean.PEPS.quantumDoubleKLatticeGauge_mul,
        TNLean.PEPS.quantumDoubleKLatticeBondPermutation_commute}
    \uses{def:peps_qd_lattice_gauge}
    The maps in (\ref{eq:qd_lattice_gauge}) satisfy $\Gamma_1=\mathrm{id}$ and $\Gamma_{uw}=\Gamma_u\circ\Gamma_w$.
    Each $\Gamma_u$ is a permutation with inverse $\Gamma_{u^{-1}}$.
    If $e\ne f$, then $V_e(s)V_f(t)=V_f(t)V_e(s)$ for every $s,t\in G$, without any flatness assumption on $x$.
\end{lemma}
\begin{proof}
    Every corner is of the form $x\mapsto u_ixu_j^{-1}$, so its successive transforms are $u_i(w_ixw_j^{-1})u_j^{-1}=(u_iw_i)x(u_jw_j)^{-1}$.
    Distinct coordinates commute in the direct product $G^E$: $\delta_e(s)\delta_f(t)=\delta_f(t)\delta_e(s)$ for $e\ne f$.
    Applying the permutation representation gives the asserted commutation.
\end{proof}

\begin{lemma}[Physical support and the four incident actions]
    \label{lem:peps_qd_lattice_bond_support}
    \lean{TNLean.PEPS.quantumDoubleKLatticeBondPermutation_right_at_left,
        TNLean.PEPS.quantumDoubleKLatticeBondPermutation_right_at_right,
        TNLean.PEPS.quantumDoubleKLatticeBondPermutation_up_at_bottom,
        TNLean.PEPS.quantumDoubleKLatticeBondPermutation_up_at_top,
        TNLean.PEPS.quantumDoubleKLatticeBondPermutation_spectator}
    \uses{def:peps_qd_lattice_gauge,def:peps_quantum_double_bond_average}
    On the edge from $v$ to $v+(1,0)$, $V_{E_v}(s)$ acts on the left and right blocks by
    \begin{align}
        (a,b,c,d)&\longmapsto(as^{-1},sb,c,d),
        & (a',b',c',d')&\longmapsto(a',b',c's^{-1},sd').
        \label{eq:qd_lattice_horizontal}
    \end{align}
    On the edge from $v$ to $v+(0,1)$, $V_{N_v}(s)$ acts on the lower and upper blocks by
    \begin{align}
        (a,b,c,d)&\longmapsto(sa,b,c,ds^{-1}),
        & (a',b',c',d')&\longmapsto(a',b's^{-1},sc',d').
        \label{eq:qd_lattice_vertical}
    \end{align}
    Every block with none of $N_v,E_v,S_v,W_v$ equal to the acted-on edge is unchanged.
\end{lemma}
\begin{proof}
    In a simple torus with periods at least three, right edges are indexed injectively by their left endpoints, up edges by their lower endpoints, and no right edge equals an up edge.
    Substituting $u=\delta_e(s)$ into (\ref{eq:qd_lattice_gauge}) therefore leaves exactly the four factors shown in (\ref{eq:qd_lattice_horizontal}) and (\ref{eq:qd_lattice_vertical}).
    At any other block all four factors equal $1$.
\end{proof}

\begin{definition}[The periodic hole holonomy]
    \label{def:peps_qd_lattice_hole}
    \lean{TNLean.PEPS.quantumDoubleKLatticePlaquetteHolonomy,
        TNLean.PEPS.quantumDoubleKLatticePlaquetteHolonomy_eq}
    \uses{def:peps_quantum_double_hole_projector,def:peps_qd_lattice_gauge}
    For a hole with lower-left block $v$, set
    \begin{align}
        B_v(x)&=b_{v+(0,1)}a_vd_{v+(1,0)}c_{v+(1,1)}.
        \label{eq:qd_lattice_hole}
    \end{align}
    This is exactly the ordered product $NW.b\,SW.a\,SE.d\,NE.c$ of Definition~\ref{def:peps_quantum_double_hole_projector} on the four indicated blocks.
\end{definition}

\begin{lemma}[Covariance of local and hole holonomies]
    \label{lem:peps_qd_lattice_holonomy_covariance}
    \lean{TNLean.PEPS.quantumDoubleKHolonomy_latticeGauge,
        TNLean.PEPS.quantumDoubleKLatticePlaquetteHolonomy_gauge,
        TNLean.PEPS.quantumDoubleKHolonomy_latticeGauge_eq_one_iff,
        TNLean.PEPS.quantumDoubleKLatticePlaquetteHolonomy_gauge_eq_one_iff}
    \uses{def:peps_qd_lattice_gauge,def:peps_qd_lattice_hole,
        def:peps_quantum_double_local_projector}
    Write $A_v(x)=a_vb_vc_vd_v$.
    For every $u\in G^E$ and every configuration $x$, one has
    \begin{align}
        A_v(\Gamma_u x)&=u_{N_v}A_v(x)u_{N_v}^{-1},
        & B_v(\Gamma_u x)&=u_{E_{v+(0,1)}}B_v(x)u_{E_{v+(0,1)}}^{-1}.
        \label{eq:qd_lattice_covariance}
    \end{align}
    Consequently $A_v(\Gamma_u x)=1$ if and only if $A_v(x)=1$, and $B_v(\Gamma_u x)=1$ if and only if $B_v(x)=1$.
\end{lemma}
\begin{proof}
    In the local product, the three adjacent pairs $u_{E_v}^{-1}u_{E_v}$, $u_{S_v}^{-1}u_{S_v}$ and $u_{W_v}^{-1}u_{W_v}$ cancel.
    For the hole product, the four successive gauge labels are $E_{v+(0,1)},N_v,E_v,N_{v+(1,0)}$ and again the initial label.
    The three intervening inverse pairs cancel, yielding (\ref{eq:qd_lattice_covariance}).
    For any $g,h\in G$, $ghg^{-1}=1$ if and only if $h=1$.
\end{proof}

\begin{definition}[Physical spins of a periodic bond coloring]
    \label{def:peps_qd_lattice_spins}
    \lean{TNLean.PEPS.quantumDoubleKLatticeSpins}
    \uses{def:peps_quantum_double_k_tensor,def:peps_qd_lattice_gauge}
    A coloring $q\in G^E$ defines the physical configuration
    \begin{align}
        F(q)_v&=\bigl(q_{N_v}q_{E_v}^{-1},\,
        q_{E_v}q_{S_v}^{-1},\,
        q_{S_v}q_{W_v}^{-1},\,
        q_{W_v}q_{N_v}^{-1}\bigr).
        \label{eq:qd_lattice_spins}
    \end{align}
\end{definition}

\begin{lemma}[Color multiplication and flatness]
    \label{lem:peps_qd_lattice_spins_flat}
    \lean{TNLean.PEPS.quantumDoubleKLatticeSpins_mul,
        TNLean.PEPS.quantumDoubleKHolonomy_latticeSpins,
        TNLean.PEPS.quantumDoubleKLatticePlaquetteHolonomy_spins}
    \uses{def:peps_qd_lattice_spins,def:peps_qd_lattice_hole,
        def:peps_qd_lattice_gauge}
    For every $u,q\in G^E$, one has $F(uq)=\Gamma_u F(q)$.
    Every coloring satisfies $A_v(F(q))=B_v(F(q))=1$ at every vertex and hole.
\end{lemma}
\begin{proof}
    Each corner satisfies $(u_iq_i)(u_jq_j)^{-1}=u_i(q_iq_j^{-1})u_j^{-1}$.
    The local product cancels cyclically.
    Substituting (\ref{eq:qd_lattice_spins}) into (\ref{eq:qd_lattice_hole}) gives the hole factors
    $q_{E_{v+(0,1)}}q_{N_v}^{-1}$, $q_{N_v}q_{E_v}^{-1}$, $q_{E_v}q_{N_{v+(1,0)}}^{-1}$ and $q_{N_{v+(1,0)}}q_{E_{v+(0,1)}}^{-1}$, whose ordered product is $1$.
\end{proof}

% These are placement and covariance statements, not a classification of
% flat configurations or of the quantum-double ground space. In particular,
% no surjectivity of F onto all flat periodic configurations is asserted.
